\begin{annexitem}{annex:p12.comb.purification}{thm:p12-comb-purification}{Purification, extension and reconstruction}
\begin{proof}
Consider first a replacement of an occurrence u in an equation
\(s=^{?}t\) by fresh b, together with the defining equation
\(b=^{?}u\). Given an old solution \(\tau\), extend it by
\(b\mapsto u\tau\). Freshness ensures that this extension leaves
all old images unchanged, and the new equations hold. Conversely,
if \(\delta\) solves the two new equations, then
\(b\delta=_Eu\delta\); replacing that equal subterm in its context
recovers \(s\delta=_Et\delta\). Splitting an equation into
\(a=^{?}s,a=^{?}t\) has the same extension/restriction proof.

For \(u=^{?}v\), identifying u with v preserves solvability in both
directions: a quotient solution expands by giving u the image of v;
an old solution has \(u\tau=_Ev\tau\), so replacing u's image by
v's image throughout preserves every other equation by congruence.
The recorded quotient substitution \(\lambda\) expresses expansion.
Iterating these arguments proves extension and restriction for the
source's alien-subterm purification.
Keeping a vanished original representative through a tautology
changes no solutions but retains its image in the pure input scope.

For the bottom-up implementation, extend an original solution by
giving each node name the value of its original subterm. Every node
equation then holds. Equal input roots have E-equal images; quotienting
their names uses the preceding identification argument. In the reverse
direction, induction on the original term tree shows that each named
node's value equals the value of its original term under the
reconstructed original-variable assignment. Equal root names thus
imply the original equations.

For coverage, extend an arbitrary original target \(\tau\) to the
purified problem as above. If a purified answer \(\sigma\) covers that
extension with factor \(\eta\), then for every original variable x
the reconstructed answer satisfies
\(x\lambda\sigma\eta=_Ex\tau\). Thus restriction preserves the
required problem-variable instantiation relation. No equality of
substitutions on fresh auxiliary variables is demanded.

The fully flattened construction visits each original node once and
creates at most one name and one defining equation per nonvariable
node. The total number of child references is at most the original
tree size. Root equations and variable equations add a linear number
of identifications. Computing the finite quotient and replacing each
reference by its representative is polynomial, even with a naive
quadratic implementation. It creates no new term nodes requiring
recursive purification. Hence it terminates and has polynomial size.
\end{proof}
\end{annexitem}

\begin{annexitem}{annex:p12.comb.soundness}{thm:p12-comb-soundness}{The ordered substitution construction}
\begin{proof}
Freshening ranges is a bijective renaming of the component's
substitutable variable alphabet, fixing all designated constants.
Extend a finite injective renaming to a permutation of the infinite
variable supply. Applying that permutation to an equation already
made \(E_i\)-equal preserves its equality. Its inverse witnesses that
the renamed answer and the original answer cover one another on
the local input variables. Constants occur in exactly the same
positions, so all LCR conditions are unchanged. Freshen identities
of unconstrained owned variables as well.

After freshening, an owned variable's image contains only its owner's
parameters and foreign shared constants. If x is shared and owned by i,
every such foreign constant y satisfies \(y<x\), by LCR and linearity
of the order. Thus each recursive definition of \(x\sigma\) uses
already defined images. There are finitely many shared variables and
then finitely many nonshared variables; each replacement is in a
finite term. This establishes well-definedness and termination.
The result's range contains parameters only, never old input variables.

Fix component i. Interpret the foreign shared constants in
\(\Gamma_i\sigma_\Pi\) as substitutable atoms temporarily and
replace each by its constructed global image. Replace parameters
by themselves. For an owned input variable x, the definition gives
\[
                   x\mu_i\sigma=x\sigma.
\]
For a foreign shared variable y, \(y\mu_i=y\), so this equality
also holds on y. It therefore holds on every variable of every
term in \(\Gamma_i\sigma_\Pi\), by induction on terms.
Because \(\mu_i\) solves this block modulo \(E_i\), substituting
the global terms for its free constants and parameters yields an
E-equality: the identities of \(E_i\) contain none of these
designated free constants, and their equational derivations remain
valid under such replacement. Hence \(\sigma\) solves the block.
Doing this for each i proves soundness for the union. Precomposing
with \(\sigma_\Pi\) and \(\lambda\) then reconstructs a solution
by the purification lemma. The displayed equality is needed only
on that block's variables, not as an equality of infinite substitutions
including other components' private variables.
\end{proof}
\end{annexitem}

\begin{annexitem}{annex:p12.comb.free-product}{thm:p12-comb-free-product}{Constructing the normalized combined algebra without a hidden support lemma}
\begin{proof}
We give all inclusion, well-definedness and universal-property steps.
For any nonempty generator set S, the quotient
\(F_i(S)=T(\Sigma_i,S)/{=_{E_i}}\) is an \(E_i\)-algebra: operations
are induced on congruence classes, and an identity holds after every
assignment of its variables to term classes. Evaluation into any
 \(E_i\)-algebra \(A\) under a map \(S\to A\) is constant on classes, since
every identity and every equational-proof step holds in A. It is
the unique homomorphism extending that generator map, by induction
on pure terms. Nontriviality keeps distinct generators distinct:
an equality of two generator classes, renamed to distinct variables,
would be the forbidden identity \(u\approx v\).

\emph{Conservative inclusion.} If \(S\subseteq S'\), fix \(z_0\in S\)
and retract S' to S by fixing S and mapping each new generator to
\(z_0\). The induced homomorphism retracts the natural inclusion
\(F_i(S)\to F_i(S')\); hence that inclusion is injective.
Suppose a nongenerator class q in \(F_i(S)\) becomes a generator
class in \(F_i(S')\). It cannot become an old generator, by
injectivity. If it becomes a new generator b, choose distinct
\(z_0,z_1\in S\) and use two retractions mapping b to \(z_0\)
and \(z_1\), respectively. They would imply
\(q=z_0=z_1\) in \(F_i(S)\), contradicting nontriviality.
Every generator set in the construction contains the infinite Z,
so the two-generator argument is always available.

\emph{Stages and the limit.} Start with distinct colour tags for the
nongenerator classes \(N_i^0\). If all nodes through stage k have
been constructed, the preceding paragraph identifies the older
pure algebra with a subalgebra of the next pure algebra. Older
nodes remain nongenerators, so the next stage retains their tags
and adds tags only for genuinely new nongenerator classes.
Tags of different colours are disjoint and never collide with named
generators. This inductively defines the displayed increasing sets.

Every pure term over \(S_i^\infty\) uses finitely many atoms. Those
atoms lie in some \(S_i^k\), so its class comes from that stage.
If it is a generator at the limit it was already that generator
at a sufficiently large stage, by conservative inclusion. Otherwise
it is a uniquely tagged i-node. Thus \(e_i\) is onto.
It is one-to-one because pure inclusions are injective, distinct
generators stay distinct, and nongenerators stay nongenerators.
For completeness, equality of two pure terms over the union is
witnessed by a finite equational derivation. Only finitely many
generators occur in that derivation; it therefore takes place in
some finite stage. This also identifies the direct limit of the
stage algebras with \(F_i(S_i^\infty)\), without assuming an
effective equality test.

Transport operations along each bijection \(e_i\). Its \(\Sigma_i\)-reduct
is thereby isomorphic to \(F_i(S_i^\infty)\), so every identity of
\(E_i\) holds in \(D\). Disjointness permits all these transported operation
families on the same carrier without conflicting interpretations.
Therefore \(D\) is an E-algebra, and its elements \(z\) in \(Z\) remain distinct.

\emph{Finite representatives and birth rank.} A node born in
\(N_i^k\) has a representative pure \(i\)-term whose generator atoms
belong to \(S_i^k\). These are z's and foreign nodes born before
k. By induction on birth rank, expand each foreign node to its
finite mixed-term representative. The resulting finite mixed term
evaluates to the node under the transported operations. Generators
are represented by themselves. Every element of D consequently has
a finite representative, and each own-colour representative uses
only lower-rank foreign nodes. No claim about a least syntactic
support, or a term over an empty generator set, has been used.

\emph{The free universal property.} Let A be any E-algebra and
\(h_0:Z\to A\) any map. Define h on nodes by increasing birth rank.
For an i-node born at stage k, take a pure representative in
\(F_i(S_i^k)\) and evaluate it in A, mapping its z-atoms by
\(h_0\) and its lower-rank foreign atoms by the already defined h.
Pure-class equality makes this value independent of the chosen
representative. More explicitly, the given generator assignment
induces a homomorphism \(F_i(S_i^k)\to A\). If a class is an old
node, its old representative is still a representative in this
algebra, so its evaluation agrees with the earlier definition.
This proves compatibility of the definitions at all stages.

For \(f\in\Sigma_i\) and a finite tuple \(d_1,\ldots,d_m\in D\),
choose a stage containing pure representatives for all
\(e_i(d_j)\), the necessary foreign atoms, and the class of their
f-combination. Evaluation of that class shows
\(h(f^D(d_1,\ldots,d_m))=f^A(h(d_1),\ldots,h(d_m))\).
This includes nullary operations, where there is no tuple.
Hence h is a homomorphism for every signature and thus for their
union. Any other homomorphism extending \(h_0\) has the same
value on each node, by its lower-rank representative and induction.
The extension is unique. D is therefore free on Z.

Finally, the ordinary combined term quotient \(T(\bigcup_i\Sigma_i,Z)/{=_E}\)
has the same free universal property, by the elementary evaluation
argument in the first paragraph. The two unique maps fixing Z are
inverse, since their composites fix generators and free extensions
are unique. Thus \(\nu(s)=\nu(t)\) exactly when \(s=_Et\).
If finitely many pure atoms are renamed injectively to distinct
generators of \(S_i^\infty\), the induced pure map is injective:
identify its image with a subset and apply the retraction argument.
Pure abstraction therefore reflects, not merely preserves,
\(E_i\)-equality. There is one empty-atom boundary case: if the two
abstracted terms are ground, there is a pure ground term \(t_0\).
The ground term quotient \(F_i(\varnothing)\) is then nonempty;
mapping every generator of \(S_i^\infty\) to \([t_0]\) retracts
its natural ground-subalgebra inclusion and proves injectivity in
this case too. If the signature has no ground terms and the local
atom list is empty, there are no terms to compare; the empty block
poses no equation. This does not invent a ground representative for
a variable-independent class such as h(z).
This proves all assertions, including collapsing and
irregular cases.
\end{proof}
\end{annexitem}

\begin{annexitem}{annex:p12.comb.projection}{thm:p12-comb-projection}{Constructing the successful local witnesses}
\begin{proof}
Set \(d_x=\nu(x\tau)\) for every purified variable. Since \(\tau\)
solves \(\Gamma\), evaluation makes each input equation equal in D.
Partition shared variables by equality of \(d_x\). This is a
genuine equivalence relation; its representatives have distinct
values. Replace each member of a class by its representative.
The resulting assignment still solves every block, since each
replacement changes an assigned value to the same element of D.

If a representative's value is an i-node, label it i. Label values
in Z with 1, as stated. For each component i, the value of a
representative labelled differently from i is either a foreign node
or a z-atom. It is therefore a generator of \(F_i(S_i^\infty)\).
The distinct values of representatives ensure that these designated
constant atoms are distinct. Choose a linear order on representatives
by nondecreasing rank and break equal-rank ties arbitrarily.

For an own shared x whose value is an i-node, choose a pure i-term
representative at its birth stage. A foreign shared constant y can
occur in that representative only if its value is one of its
foreign generator atoms. Such an atom has strictly smaller rank
than \(d_x\); hence \(y<x\). If \(d_x\in Z\), use a single
parameter representing that generator; it contains no foreign
shared constant. Thus every own shared image satisfies LCR.
Nonshared variables have no image restrictions by the source's
placement after all shared representatives. For an unconstrained
own representative absent from \(\Gamma_i\), use the same birth-stage
image; adding it cannot affect the equations.

Only finitely many generator atoms occur in these chosen images
and in images for the nonshared variables. Include every designated
foreign shared constant in this finite atom list. Name its value by
that constant, and name all other atoms by distinct fresh parameters.
This gives an injective atom map \(\alpha_i\) from the local names
to \(S_i^\infty\). It does not identify a new parameter with a
designated constant. The local substitution \(\theta_i\) maps an
owned variable to the abstracted pure representative and fixes
foreign constants.

Under \(\alpha_i\), the class of \(x\theta_i\) is
\(e_i(d_x)\) for every variable in the i-block: this is the chosen
representative for own variables, and a generator for foreign ones.
Consequently, for a pure term s in that block, induction on s
identifies the class of \(s\theta_i\) under \(\alpha_i\) with
\(e_i(s^D(d))\). Because the input equation \(s=^{?}t\) holds
under d in D, these classes are equal. The injective pure atom
map reflects equality by the preceding theorem, so
\(s\theta_i=_{E_i}t\theta_i\) in the local term language.
This proves local solvability as well as all the restrictions.
The construction uses the proved combined algebra rather than
the potentially misleading roots of the written images \(x\tau\).
\end{proof}
\end{annexitem}

\begin{annexitem}{annex:p12.comb.completeness}{thm:p12-comb-completeness}{Actual factorization through a chosen tuple of complete answers}
\begin{proof}
By purification, extend any original target solution to a target
\(\tau\) of the pure problem. The projection lemma supplies a branch
and local LCR targets \(\theta_i\). Since \(C_i\) is complete,
choose \(\mu_i\in C_i\) and a constant-fixing factor \(\rho_i\)
such that, on every owned local variable x,
\[
                  x\mu_i\rho_i=_{E_i}x\theta_i.
\]
Range freshening does not affect existence of the factors. Tautological
constraints include every unconstrained owned representative in the
local input scope, so its target is covered by the same local factor.
The
parameters \(P_i\) of different answers are disjoint.

Give each parameter p in \(P_i\) a value in D as follows. Evaluate
the pure term \(p\rho_i\) in the i-reduct, giving each local target
parameter the generator value specified by \(\alpha_i\), and each
designated foreign constant y the value \(d_y\).
If \(\rho_i\) uses extra unconstrained parameters, give them arbitrary
generator values as well. The finitely many factor equalities only
involve finitely many such atoms. Choose a finite mixed-term
representative of the resulting D-value for each p; such a representative
exists by the free-product theorem. Define the global factor \(\eta\)
on all these p's by those mixed terms, and leave other variables fixed.
Disjointness of \(P_i\) makes this one finite substitution, with no
conflicting assignments. It need not preserve component LCR under
instantiation, nor does completeness require that: it is a factor
for a combined target, not a second component answer.

Let \(\sigma\) be the ordered merge. We prove
\(\nu(x\sigma\eta)=d_x\) by induction on the shared order.
For x owned by i, the merged image is \(x\mu_i\) with each
foreign constant y replaced by \(y\sigma\). Such y precedes x,
so its value after \(\eta\) is \(d_y\) by induction. The parameters
of \(x\mu_i\) have exactly the values just assigned through
\(\rho_i\). The local factor equality therefore evaluates to
\[
 \nu(x\sigma\eta)
 = (x\mu_i\rho_i)^{F_i(S_i^\infty)}\alpha_i
 = (x\theta_i)^{F_i(S_i^\infty)}\alpha_i
 = d_x,
\]
where the displayed pure classes are identified with their D-values
through \(e_i^{-1}\). The least shared variable has no foreign
constant dependencies, so the induction starts. For nonshared
variables the same computation applies after all shared variables.
All ranges in the merge contain parameters only; substituting
\(\eta\) therefore effects exactly these replacements, even if
the target's written terms contain original problem-variable names.

Equality in D is equivalent to E-equality of finite mixed terms.
Thus \(x\sigma\eta=_Ex\tau\) on the purified input variables.
For a shared variable not chosen as representative, its target
value equals its representative's by the definition of \(\Pi\),
so \(\sigma_\Pi\sigma\eta\) covers it as well. Finally the
purification reconstruction covers the original target on all
original variables. This proves completeness of the union of all
branches and tuples.

There are finitely many branches on a finite shared-variable set.
If each \(C_i\) is finite, each branch has a finite Cartesian product,
and the union is finite. A finite complete set in a preorder has
a minimal complete subset: discard duplicate equivalence classes
and any class strictly above another retained class in the
more-general ordering, until only minimal classes remain.
Finite descent terminates and each discarded answer is covered
by a retained answer. This last existence argument does not
claim an effective global comparison procedure when none is supplied.
Hence the combined theory is at most finitary, with no claim that
the raw generated union is minimal.
\end{proof}
\end{annexitem}

\begin{annexitem}{annex:p12.comb.ac-lcr}{thm:p12-comb-ac-lcr}{Fixed-atom persistence and the nested AC construction}
\begin{proof}
Every pure AC term is represented by its nonempty multiset of variable
and designated constant atoms, by Theorem~\ref{thm:p12-ac-multisets}.
For a constant c fixed by a factor substitution \(\rho\), its
multiplicity in an image \(x\mu\rho\) is its multiplicity in
\(x\mu\), plus a sum of nonnegative contributions from the
parameter images. Consequently an occurrence of c in \(x\mu\)
cannot disappear in any AC-instance of \(\mu\). No cancellation
of opposite signs or removal by a unit is available in this theory.

The flat algorithm returns a finite complete family C on the full
input variable scope, including tautological/unconstrained variables.
Every retained answer is an AC unifier and satisfies every LCR
occurrence test, so the filtered family is sound. Let \(\tau\) be
any permitted AC-LCR target. There is \(\mu\in C\) and a
constant-fixing factor \(\rho\) with
\(x\mu\rho=_{AC}x\tau\) for every input variable x.
If \(\mu\) contained a forbidden c in x's image, persistence
would force c in \(x\tau\), contradicting permission. Thus the
covering answer survives all tests. The filtered family is complete.
All tests traverse finitely many finite images of a finite family,
so the algorithm terminates. It also decides the pure LCR problem
by testing whether this finite family is empty.

For a finite nested input, use one disjoint component signature
\(\{f_j\}\) for each AC symbol \(f_j\), with its two AC axioms.
Their signatures are disjoint. Put all uninterpreted operations and
free constants in one separate empty-theory component. Purification
may leave repeated/nested occurrences of a component's own binary
AC operation; flattening those occurrences produces exactly the
finite pure \(f_j\)/variable/foreign-constant problem accepted by the
flat solver, not a problem with other positive-arity symbols in
that component. The present paragraph and the syntactic LCR proof
supply terminating finite complete component procedures.
The decision and completeness theorems therefore apply to the
finite branch search and its actual ordered merges. They prove
termination, soundness and coverage for arbitrary finite nesting
and for any finite number of such AC symbols. Positivity of each
flat component image is preserved when a foreign constant is
replaced by a nonempty finite global term. No step assigns an
empty image to a selected AC parameter or inserts an ACU unit.
This is a genuine bridge from the already proved flat construction,
not an appeal to the mere name of a combination theorem.
\end{proof}
\end{annexitem}

\begin{annexitem}{annex:p12.comb.decision}{thm:p12-comb-decision}{Finite branch search, NP certificates and syntactic LCR}
\begin{proof}
For m shared variables, a partition can be encoded by m block indices,
a label by m indices in \(\{1,\ldots,n\}\), and an order by a
permutation of at most m representatives. All choices form a finite
effectively enumerable set; there are at most
\(B_m n^m m!\) choices, where \(B_m\) is the number of partitions.
Purification has polynomial size by its proof. Every induced LCR
problem is finite and effectively constructed from that data.
Run its terminating correct component decision procedures. Accept
exactly when a branch has all components solvable. Soundness and
projection prove the equivalence with original solvability. Finite
enumeration and terminating component calls prove termination.

If each local problem is in NP for the fixed theories, guess the
partition, labels and order in polynomial space/time to write their
encoding, then guess and check the component NP certificates.
There are fixed finitely many components; their induced input
sizes and certificate bounds are polynomial in the original size.
The conjunction of their checks and the branch verification are
therefore polynomial. Soundness/completeness give an NP decision
procedure; the tower is used only to prove existence of a branch,
not computed as part of the certificate.

For finite complete-set generation, compute the local finite sets
in every branch, freshen their ranges, enumerate each finite
Cartesian product, and perform the finite ordered merge. Each
stage terminates under the stated local hypotheses. Duplicate
syntactically identical reconstructed answers can be removed.
Completeness is the preceding theorem; no generality comparison
is needed to return a complete, possibly redundant family.

For the free component, use the finite-tree syntactic unification
interface with designated constants fixed. Failure of ordinary
unification means there is no LCR solution. Otherwise take its
MGU \(\mu\) on the local variables and freshen variable ranges.
Every syntactic solution is \(\mu\rho\) on that scope, where
\(\rho\) fixes the constants. If a forbidden constant c occurs
in \(x\mu\), that same occurrence survives every such
instantiation: substitutions can replace variable leaves, not
delete a constant leaf or its function-symbol ancestors. Thus
no solution can satisfy that restriction. If no forbidden
occurrence exists, the MGU itself is permitted and covers every
permitted solution. This proves the proposed LCR singleton/failure
test, including unconstrained variables after freshening.

Finally, the free symbols of a finite general input form a finite
signature disjoint from the theory signatures. The empty theory
on it is nontrivial and has exactly syntactic equality. Include
it as one more component and apply the already proved theorem.
Every mixed free-symbol term is purified as an alien term at a
different component, so none is incorrectly treated as an atom
inside a theory equation without its defining constraint.

Under terminating finite complete-set procedures the component
decision interfaces are available by testing emptiness of their
returned finite sets. We can consequently compare any two generated
answers \(\sigma,\tau\) effectively on the finite original scope X.
Rename the variables in \(\sigma\)'s images to a disjoint fresh
left alphabet. Replace every variable in \(\tau\)'s images by
its own distinct fresh constant \(c_z\). The test equations are
\(x\sigma'=^{?}_E(x\tau)\kappa\), one for each \(x\in X\),
where \(\kappa:z\mapsto c_z\). The only substitutable variables
in these equations are in the left alphabet. The new constants
form an additional free signature, or are added to the existing
free component, so the proved disjoint combination decision method
applies without changing any \(E_i\).

If \(\sigma\) covers \(\tau\) by a factor \(\eta\), rename its
domain to the left alphabet and replace target-variable occurrences
in its images by their \(c_z\)'s. This gives a solution of the test:
an equational derivation remains valid under variable-to-constant
replacement because these are fresh symbols absent from E.
Conversely, from a solution of the test replace every \(c_z\) by \(z\)
and undo the left-alphabet renaming. Equational replacement now
gives a finite factor with \(x\sigma\eta=_Ex\tau\) for all x.
Other fresh image variables remain parameters; no idempotence of
the factor is required. The test is therefore exactly coverage,
not merely equality of the two written answers.

Run these terminating tests on all ordered pairs of the finite
complete output. Mutual coverage forms equivalence classes. The
induced preorder on those classes is a finite partial order; retain
one answer from each class having no strictly more-general class.
Every discarded class has a finite descending chain to a retained
class, so coverage is preserved by transitivity. Distinct retained
classes are incomparable. The resulting finite family is minimal
complete, and the whole comparison/pruning procedure terminates.
Without the terminating local algorithms the preceding theorem
asserts finite minimal-set existence only, not this effective step.
\end{proof}
\end{annexitem}
